\documentclass[11pt,a4paper]{article}
\usepackage{../shared/hanzocover}
\usepackage[utf8]{inputenc}
\usepackage[T1]{fontenc}
% Shared preamble for the compute-market series.
% One style, one vocabulary, inputted by every paper so the series reads as one voice.
\usepackage[margin=1in]{geometry}
\usepackage{booktabs}
\usepackage{amsmath}
\usepackage{amssymb}
\usepackage{amsthm}
\usepackage{graphicx}
\usepackage{xcolor}
\usepackage{hyperref}
\hypersetup{colorlinks=true, linkcolor=blue, urlcolor=blue, citecolor=blue}
\usepackage{enumitem}
\usepackage{listings}
\usepackage{array}
\usepackage{tabularx}
\usepackage{siunitx}
\sisetup{group-separator={,}, group-minimum-digits=4}

\definecolor{kw}{rgb}{0.13,0.29,0.53}
\definecolor{cmt}{rgb}{0.40,0.40,0.40}
\definecolor{str}{rgb}{0.16,0.50,0.16}
\lstset{
  basicstyle=\ttfamily\footnotesize,
  keywordstyle=\color{kw}\bfseries,
  commentstyle=\color{cmt}\itshape,
  stringstyle=\color{str},
  showstringspaces=false,
  breaklines=true,
  columns=fullflexible,
  frame=single,
  framesep=4pt,
}

\newtheorem{definition}{Definition}
\newtheorem{proposition}{Proposition}
\newtheorem{remark}{Remark}

\newcommand{\code}[1]{\texttt{#1}}
\newcommand{\repo}[1]{\texttt{#1}}

% Series vocabulary — used identically in every paper.
\newcommand{\job}{\ensuremath{J}}
\newcommand{\bundle}{\ensuremath{R}}
\newcommand{\cost}{\ensuremath{\mathcal{C}}}

\tolerance=800
\emergencystretch=3em


\title{Threshold Settlement:\\Escrow and Payout Without a Key}

\author{
Zach Kelling\thanks{Corresponding author: research@lux.network} \\
\textit{Lux Partners} \\
\textit{A subsidiary of Hanzo Industries Inc} \\
\texttt{https://lux.network}
}

\date{2026}

\begin{document}
\hanzocoverpage

\begin{abstract}
Every compute market eventually holds someone else's money. A job is posted,
funds are escrowed, work is performed, a verifier decides, and a payout is
issued. The question nobody enjoys asking is who holds the key that releases the
escrow.

The answer should be nobody. We describe a settlement architecture in which the
private key authorising payouts never exists --- not in memory, not on disk, not
during signing --- and in which the operator set can change without the key ever
being reconstructed. We give the protocol table we run in production across
twenty or more chains, including two post-quantum lattice families and a
composition requiring both, and report measured signing and key-generation
latencies.

We then describe the interface boundary that makes this composable with a market
rather than entangled with it: an assembler that turns a payment specification
into an unsigned transaction and a set of digests, and a finaliser that turns
digests plus signatures back into a broadcastable transaction. The market decides
what to pay. The threshold cluster decides whether to sign. Neither needs to
understand the other.
\end{abstract}

\tableofcontents
\newpage

\section{The custody question}

Escrow in a decentralized market is usually solved one of three ways, and two of
them are bad.

A single operator key is simple and is a single point of theft, a single point
of coercion and a single point of loss. A smart contract on the settlement chain
is genuinely good where the asset is native to that chain, and useless where it
is not --- which is the common case as soon as providers want paying in something
that lives elsewhere. A multi-signature scheme is better than one key but
publishes the entire policy on chain, costs linearly in signers, and is
unavailable or awkward on many chains.

Threshold signing gives the properties of multi-signature without the
disclosure, the cost or the chain dependence: $t$ of $n$ parties cooperate to
produce a signature that is indistinguishable from a single-key signature, and
the full key is never assembled at any point in the protocol.

\begin{definition}[No-reconstruction property]
A threshold signing scheme has the no-reconstruction property if there exists no
point in key generation, signing, or resharing at which any single party holds
material sufficient to recover the private key, and no protocol transcript from
which fewer than $t$ colluding parties can recover it.
\end{definition}

This is stronger than secret sharing with reassembly at signing time, which is
what several deployed systems actually do and which reintroduces the single
point of theft for the duration of every signature.

\section{The protocol table}

\begin{table}[h]
\centering
\begin{tabular}{lll}
\toprule
Protocol & Curve or lattice & Settlement reach \\
\midrule
CGGMP21 & secp256k1 & Bitcoin, Ethereum, all EVM chains, XRPL \\
FROST & Ed25519 & Solana, TON, Polkadot, Cardano, Substrate \\
BLS & BLS12-381 & aggregate attestation, beacon signatures \\
SR25519 & Ristretto255 & Substrate native \\
Pulsar & M-LWE (threshold ML-DSA-65) & post-quantum, anchors and custody \\
Corona & R-LWE & post-quantum, orthogonal lattice family \\
Double lattice & Pulsar $\oplus$ Corona & both must verify \\
\bottomrule
\end{tabular}
\caption{Threshold protocols in production, with the settlement surfaces each
reaches.}
\end{table}

Three properties of this table matter for a market.

\paragraph{Non-interactive signing.} CGGMP21 requires an interactive key
generation once, after which signing is non-interactive. For a market issuing
many payouts this is the difference between a settlement path that scales and
one that turns into a coordination problem at load. It also produces signatures
bit-identical to any standard secp256k1 signature, so nothing downstream needs
to know a threshold scheme was involved.

\paragraph{Dynamic resharing.} Parties can be added or removed without downtime
and without reconstructing the key. In a permissionless market the operator set
churns; a custody design that requires a ceremony every time it changes will not
survive contact with reality.

\paragraph{Post-quantum options with a hedge.} A custody key is the longest-lived,
highest-value secret in the system, which makes it the worst possible thing to
leave on an elliptic curve. Classical threshold ECDSA is $t$-of-$n$ against a
classical adversary and $0$-of-$n$ against Shor. The lattice families provide a
post-quantum threshold; the double-lattice composition, which concatenates a
Pulsar and a Corona signature and requires both to verify, hedges against a break
in one lattice problem rather than in lattices generally.

\section{Measured behaviour}

Signing completes in under \SI{25}{\milli\second} and key generation in
\SIrange{12}{82}{\milli\second}, with Byzantine fault tolerance up to $t-1$
malicious parties. For a settlement path these are not the binding constraint ---
block times dominate by two to three orders of magnitude --- which is the
comfortable position to be in.

\section{The interface boundary}

The design decision that makes this composable is refusing to let the signer know
what a job is.

\begin{lstlisting}
PreSign(spec)            -> (unsigned transaction, per-input digests)
Finalize(unsigned, sigs) -> (raw transaction, txid)
\end{lstlisting}

The assembler takes a payment specification --- source, destination, amount, fee
policy --- and returns a wire-correct unsigned transaction together with the
digests that must be signed, computed under the correct sighash rules for the
chain. The finaliser takes those digests' signatures and produces a
broadcastable transaction. Between the two calls sits whatever authorisation
process the operator wants, and the assembler neither knows nor cares what it is.

This yields a clean separation for a market:

\begin{itemize}[leftmargin=1.4em]
\item The \emph{allocator} decides who should be paid and how much. It has no
      access to key material and cannot move funds.
\item The \emph{verifier} decides whether the work was done. It produces an
      attestation, not a signature.
\item The \emph{threshold cluster} decides whether to sign a specific
      transaction against a specific attestation policy. It has no opinion about
      pricing.
\end{itemize}

Each can be audited, replaced, or run by a different party. A compromise of the
allocator produces mispriced jobs, not stolen funds. A compromise of the verifier
produces incorrect attestations, which the signing policy can require multiple of.

\section{Per-job escrow}

We recommend deriving a distinct escrow address per job rather than pooling.

\begin{equation}
\text{escrow}(\job) \;=\; \mathrm{derive}\big(\text{master},\; \text{path}(\job)\big),
\end{equation}
with a deterministic path incorporating the job identifier. The properties this
buys are worth the small extra bookkeeping. Funds for distinct jobs are isolated,
so a policy failure on one job cannot drain another. The escrow address is
computable by the buyer before funding, so the buyer can verify where its money is
going without trusting an API response. Settlement is a single threshold
signature over a transaction whose specification the verifier has already
validated. And an unclaimed escrow is visibly unclaimed on chain, which makes
refund logic auditable rather than a database state.

The derivation should use a hardened path structure with one seed and many paths
rather than many seeds. Maintaining $n$ independent seeds multiplies custody risk
with no compensating benefit; independent authorisation boundaries are better
enforced at the access-control layer, where they can be scoped and revoked, than
by seed proliferation.

\section{What a market has to decide}

Threshold signing removes the key-holder question but does not answer the policy
question, and the policy is where the real design work is.

\paragraph{What does the cluster require before signing?} A verifier attestation
is the minimum. Two independent attestations, or an attestation plus an elapsed
challenge window, are stronger. The window matters: it converts an undetected
verifier compromise into a race the honest party can win.

\paragraph{What is the refund path?} A job whose deadline passes without a valid
result must return funds. This should be a signing policy the cluster can execute
against an on-chain timestamp, not a manual operation.

\paragraph{What is the dispute path?} If the buyer and verifier disagree, whose
attestation governs? We suggest the answer is neither: escalate to a higher
verification tier and let the result decide, with the loser paying for it. That
makes frivolous disputes self-limiting.

\paragraph{Who runs the parties?} The security of a $t$-of-$n$ scheme is the
security of the least correlated $t$ parties. Parties operated by one entity, on
one cloud, in one region, under one jurisdiction are not $n$ parties. This is the
same diversity problem replication has, and it deserves the same explicit
treatment rather than being assumed away.

\section{Availability}

The threshold signing engine, the multi-chain adapters, the transaction assembler
and the key-management integration are open source under permissive licence. The
signing engine can be adopted without the assembler, and the assembler without
the market --- the boundary described above is real, not aspirational, and each
piece is separately useful.

\end{document}
